\scanpage{34}
% archive sheet number: 33
% At far right a sliver of the leaf beneath bears a partial A-shaped mark.
Soit $G$ un \inserted{schéma en} groupe affine sur $A$, $A$-simple

\[
\begin{aligned}
\operatorname{Aff}_A(G)&\qquad\text{son algèbre affine}\\
\mathcal U_A(G)&\qquad\text{son hyperalgèbre}
\end{aligned}
\]

Donc $\deleted{\unsure{\{|\ell}}$ $\operatorname{Aff}_A(G)$ est une $A$-algèbre commutative augmentée, munie d'un \unsure{hom|homom.} diagonal multiplicatif, associatif, compatible avec les augmentations

$U_A(G)$ est une $A$-algèbre en général non comm., munie d'un hom diagonal multiplicatif, associatif, comm. compatible avec les augmentations

Il y a un produit scalaire

\[
\langle f,X\rangle=\langle\deleted{\unsure{f}}X,\deleted{\unsure{X}}f\rangle
\]

entre $\operatorname{Aff}_A(G)$ et $U_A(G)$, ce dernier s'identifiant aux formes $A$-linéaires sur $\operatorname{Aff}_A(G)$ qui s'annulent sur une puissance de l'idéal d'augmentation

$U_A(G)$ opère à g. et à droite sur $\operatorname{Aff}_A(G)$ [qui est un bimodule sur $U_A(G)$]

\[
\left\{\begin{aligned}
\langle Y,X*f\rangle&=\langle Y*X,f\rangle=\langle f,Y*X\rangle=\langle f*Y,X\rangle\\
\langle\varepsilon,X*f\rangle&=\langle\varepsilon,f*X\rangle=\langle X,f\rangle
\end{aligned}\right.
\]
